We use a four-step proof spine.
The purpose is to keep the main argument transparent for mathematical audit.

\subsection*{Main proof spine}
\begin{enumerate}
  \item \textbf{Define balance.}
  Fix $\mathcal{A}$, define $\Phi_H$, and define
  \[
  \mathcal{B}=\{s\in\mathcal{A}_{\mathrm{strip}}:\Phi_H(s)=0\}.
  \]

  \item \textbf{Target A (correspondence).}
  Prove that nontrivial zeros of $\zeta$ correspond exactly to points in $\mathcal{B}$
  (\cref{thm:target-correspondence}).

  \item \textbf{Target B (critical-line exclusion).}
  Prove that all points of $\mathcal{B}$ satisfy $\RePart(s)=\tfrac12$
  (\cref{thm:stage6-master-exclusion-target}).

  \item \textbf{Deduce RH.}
  Combine Steps 2--3 to obtain the conditional RH reduction
  (\cref{prop:rh-reduction-main-spine}).
\end{enumerate}

\subsection*{Status}
Step~1 is definitional.
Steps~2 and~3 are open theorem targets in this manuscript version.
Therefore Step~4 is conditional, and RH is not claimed as proved here.

\subsection*{Where technical machinery now lives}
Detailed transfer frameworks, defect-function constructions, and Stage~6 subroutines are retained as technical supporting material (appendices/notes), not as part of the main proof spine. In particular, the prior long OP/AC cross-cutting catalog is no longer in the main proof architecture; only a short dependency acknowledgment remains in manuscript-level summaries.
